% =====================================================================
\section{Results}
\label{sec:results}
% =====================================================================
% RULES FOR THIS SECTION
%  - No conclusions yet. Every number must be generated from repository
%    outputs (never typed by hand from README prose).
%  - Primary numeric sources:
%      outputs/a1_a5_behavior_analysis/            (Layer 1, reference-free)
%      outputs/evaluation_protocol_v2/<model>/<exp>/A*/A*_metrics.json  (Layer 2)
%      outputs/semantic_gt_review/human_review.csv (Layer 3; read-only)
%  - Use "LabelMe-matched" / "LabelMe-unmatched". Never call an unmatched
%    detection a "false positive" or "non-palm" unless the semantic review
%    supports it.
%  - Sub-5,747 runs: use raw counts; *_pct fields in their metrics JSON are
%    divided by 5,747 (known evaluator quirk).

\TODO{Write results prose only after all tables and figures below are generated. Report observations, not explanations; explanations belong in \cref{sec:discussion}.}

% ---------------------------------------------------------------------
\subsection{Overall Verification Behavior}
\label{sec:results:overall}
% ---------------------------------------------------------------------
% Answers: context for RQ1--RQ3.
% Content: for each model with full-scale runs, the R/U/Ur distribution,
% coverage and abstention under every condition it was run on; the
% always-Reliable baseline as a reference row for Layer-2 metrics.

% TODO: AUTO-GENERATE FROM RESULTS
% TODO: AUTO-GENERATE FROM RESULTS
% Source: outputs/a1_a5_behavior_analysis/table_b_decision_distribution.csv
%         (6 complete models), extended with every scaling-campaign run and
%         MiniCPM/Molmo2 A2--A5 once complete. Cross-check counts against the
%         verification_label column of Protocol-v2 evaluation CSVs.
% Columns: counts and % of N=5,747; coverage = (R+Ur)/N; abstention = U/N.
\begin{table}[t]
  \centering
  \small
  \caption{Decision distribution (\R{}/\U{}/\Ur{}) per model and input condition, $N=5{,}747$. Reference-free. \AUTOGEN{caption details}}
  \label{tab:decision_distribution}
  \begin{tabular}{@{}l l r r r r r@{}}
    \toprule
    \textbf{Model} & \textbf{Cond.} & \textbf{\R{}} & \textbf{\U{}} & \textbf{\Ur{}} & \textbf{Coverage} & \textbf{Abstention} \\
    \midrule
    \AUTOGENROW{7}{one row per model $\times$ condition} \\
    \bottomrule
  \end{tabular}
\end{table}


% TODO: AUTO-GENERATE FROM RESULTS
% TODO: AUTO-GENERATE FROM RESULTS
% Source: outputs/evaluation_protocol_v2/<model>/<experiment_id>/A*/A*_metrics.json
%   (fields: accuracy, precision, recall, f1, specificity, balanced_accuracy,
%    true_positive, false_positive, false_negative, true_negative,
%    evaluation_protocol == "v2").
% Include an "always-Reliable" baseline row (acc 0.889, spec 0, F1 from prior)
%   computed by the same script, not typed by hand.
% Use only full-cohort (N=5,747) runs. Do not mix @1000 runs into this table.
\begin{table}[t]
  \centering
  \small
  \caption{Protocol-v2 annotation-alignment metrics, model $\times$ condition ($N=5{,}747$; \U{} excluded). ``Spec'' is alignment specificity on \unmatched{} detections, not semantic non-palm specificity. \AUTOGEN{caption details}}
  \label{tab:alignment_metrics}
  \begin{tabular}{@{}l l r r r r r r r@{}}
    \toprule
    \textbf{Model} & \textbf{Cond.} & \textbf{Cov.} & \textbf{Acc} & \textbf{Prec} & \textbf{Sens} & \textbf{Spec} & \textbf{BalAcc} & \textbf{F1} \\
    \midrule
    Always-\R{} baseline & -- & \multicolumn{7}{c}{\AUTOGEN{baseline row}} \\
    \midrule
    \AUTOGENROW{9}{one row per model $\times$ condition} \\
    \bottomrule
  \end{tabular}
\end{table}


\TODO{Describe the decision distributions and alignment metrics (\cref{tab:decision_distribution,tab:alignment_metrics}) without interpretation. Report coverage next to every accuracy/F1 value.}

% ---------------------------------------------------------------------
\subsection{Effect of Visual Context}
\label{sec:results:context}
% ---------------------------------------------------------------------
% Answers: RQ1. Within-model comparison of A1--A5.
% Contrasts to report separately:
%   metadata-only (same image):  A1 vs A2 (confidence), A2 vs A3 (geometry)
%   image-only (same text):      A2 vs A4 (add crop panel), A2 vs A5 (crop only),
%                                A4 vs A5 (remove full-patch panel)
%   A1 vs A4/A5 changes image AND text — label as such.

% TODO: AUTO-GENERATE FROM RESULTS
\begin{figure}[t]
  \centering
  \FIGPLACEHOLDER{Context-ablation plot: per model, stacked R/U/Ur shares for A1--A5 (source: \texttt{outputs/a1\_a5\_behavior\_analysis/figures/decision\_distribution\_by\_condition.png}; regenerate in thesis style).}
  \caption{Decision distribution by input condition for each model. \AUTOGEN{caption details}}
  \label{fig:context_distribution}
\end{figure}

% TODO: AUTO-GENERATE FROM RESULTS
\begin{figure}[t]
  \centering
  \FIGPLACEHOLDER{Context-ablation plot: alignment specificity, sensitivity, balanced accuracy and coverage vs.\ condition, one line per model, with image-clustered bootstrap CIs.}
  \caption{Protocol-v2 alignment metrics by input condition. These are annotation-alignment metrics, not semantic accuracy. \AUTOGEN{caption details}}
  \label{fig:context_alignment}
\end{figure}

% TODO: AUTO-GENERATE FROM RESULTS
% TODO: AUTO-GENERATE FROM RESULTS
% Source: outputs/a1_a5_behavior_analysis/pairwise_transitions.csv
%         (all condition pairs) and table_c_transitions_a1_reference.csv.
% Report separately:
%   metadata-only contrasts (same image): A1->A2, A2->A3
%   image-only contrasts (same text):     A2->A4, A2->A5, A4->A5
\begin{table}[t]
  \centering
  \small
  \caption{Per-detection decision changes between conditions: share changed, toward rejection, toward acceptance, net. Direction labels describe the shift along \R{}$\rightarrow$\U{}$\rightarrow$\Ur{}; neither direction is better. \AUTOGEN{caption details}}
  \label{tab:transitions}
  \begin{tabular}{@{}l l l r r r r@{}}
    \toprule
    \textbf{Model} & \textbf{Contrast} & \textbf{Type} & \textbf{Changed} & \textbf{$\rightarrow$Rej.} & \textbf{$\rightarrow$Acc.} & \textbf{Net} \\
    \midrule
    \AUTOGENROW{7}{rows for A1$\rightarrow$A2, A2$\rightarrow$A3 (metadata-only); A2$\rightarrow$A4, A2$\rightarrow$A5, A4$\rightarrow$A5 (image-only)} \\
    \bottomrule
  \end{tabular}
\end{table}


% TODO: AUTO-GENERATE FROM RESULTS
\begin{figure}[t]
  \centering
  \FIGPLACEHOLDER{Per-model $3\times3$ transition heatmaps for the image-only contrasts A2$\rightarrow$A4 and A2$\rightarrow$A5 (source script: \texttt{scripts/visualization/plot\_a1\_a5\_behavior.py}; existing figure uses A1 reference).}
  \caption{Per-detection decision transitions between conditions. \AUTOGEN{caption details}}
  \label{fig:transitions}
\end{figure}

\TODO{Report metadata-only and image-only contrasts separately. Note any condition in which a model's behavior changes sharply.}

% ---------------------------------------------------------------------
\subsection{Cross-Model and Scaling Analysis}
\label{sec:results:scaling}
% ---------------------------------------------------------------------
% Answers: RQ2.
% Valid comparisons:
%   Across families: models with complete A1--A5 @5,747 on the same
%     cohort (currently Qwen2.5-VL-7B, Qwen3-VL-8B, GLM-4.6V-Flash,
%     Phi-4-multimodal, InternVL3.5-8B-HF, Gemma 4 12B; + MiniCPM-V-4.5,
%     Molmo2-8B once A2--A5 finish).
%   Within family (scale ladders, same family adapter/env):
%     Qwen3-VL 2B / 4B / 8B (/ 32B if probe passes)
%     InternVL3.5-HF 2B / 4B / 8B / 14B
%     Qwen2.5-VL 3B / 7B (/ 32B if probe passes)
%     Ministral 3 3B / 8B / 14B (only if probes pass and runs complete)
% Caveats to state with every scaling plot: anchors (8B/7B) were run
% earlier; Qwen2.5-VL-7B revision unrecorded; attention backend resolved
% to SDPA for anchors per campaign doc; hardware differences (H200 vs L40S).
% Do NOT compare across families by parameter count alone.

% TODO: AUTO-GENERATE FROM RESULTS
\begin{figure}[t]
  \centering
  \FIGPLACEHOLDER{Scaling plot: per family (Qwen3-VL, InternVL3.5-HF, Qwen2.5-VL), x = parameter count (log), y = R/U/Ur shares and coverage, one panel per condition.}
  \caption{Decision behavior versus model size within each family. \AUTOGEN{caption details}}
  \label{fig:scaling_behavior}
\end{figure}

% TODO: AUTO-GENERATE FROM RESULTS
\begin{figure}[t]
  \centering
  \FIGPLACEHOLDER{Scaling plot: context sensitivity versus size, e.g.\ share of decisions changed A2$\rightarrow$A5 and A2$\rightarrow$A4 per model size, within family.}
  \caption{Sensitivity to input condition versus model size. \AUTOGEN{caption details}}
  \label{fig:scaling_sensitivity}
\end{figure}

% TODO: AUTO-GENERATE FROM RESULTS
% TODO: AUTO-GENERATE FROM RESULTS
% Not yet computed anywhere in the repository. Proposed computation
% (reference-free, from stored decisions):
%   - pairwise Cohen's kappa (3-class) between models within each condition;
%   - sign agreement of condition effects: for each contrast (e.g. A2->A5),
%     does each model move toward rejection, toward acceptance, or neither?
\begin{table}[t]
  \centering
  \small
  \caption{Consistency of condition effects across models: direction of the net shift for each contrast, and mean pairwise inter-model agreement per condition. \AUTOGEN{caption details}}
  \label{tab:cross_model_agreement}
  \begin{tabular}{@{}l c c c c c@{}}
    \toprule
    \textbf{Model} & \textbf{A1$\rightarrow$A2} & \textbf{A2$\rightarrow$A3} & \textbf{A2$\rightarrow$A4} & \textbf{A2$\rightarrow$A5} & \textbf{A4$\rightarrow$A5} \\
    \midrule
    \AUTOGENROW{6}{one row per model; cells = direction / magnitude of net shift} \\
    \midrule
    Mean pairwise $\kappa$ & \multicolumn{5}{c}{\AUTOGEN{per condition}} \\
    \bottomrule
  \end{tabular}
\end{table}


\TODO{Describe whether the direction of condition effects (\cref{sec:results:context}) agrees across families and sizes. Blocked until scaling runs finish (see \texttt{ANALYSIS\_PLAN.md}).}

% ---------------------------------------------------------------------
\subsection{Uncertainty and Abstention Behavior}
\label{sec:results:abstention}
% ---------------------------------------------------------------------
% Answers: RQ3.
% Content: R/U/Ur distributions (Table tab:decision_distribution), coverage
% vs condition, where Uncertain mass goes when the condition changes
% (transition matrices), and per-detection stability.

% TODO: AUTO-GENERATE FROM RESULTS
\begin{figure}[t]
  \centering
  \FIGPLACEHOLDER{Ternary or stacked-bar plot of (R, U, Ur) shares for every model $\times$ condition cell, models grouped by family.}
  \caption{Use of the three verdicts across models and conditions. \AUTOGEN{caption details}}
  \label{fig:ruur_simplex}
\end{figure}

% TODO: AUTO-GENERATE FROM RESULTS
\begin{figure}[t]
  \centering
  \FIGPLACEHOLDER{Coverage (= (R+Ur)/N) vs.\ alignment balanced accuracy on decided samples, one point per model $\times$ condition.}
  \caption{Coverage versus alignment balanced accuracy. Points with low coverage are evaluated on few samples. \AUTOGEN{caption details}}
  \label{fig:coverage}
\end{figure}

\TODO{Describe verdict usage per model, without labeling a model ``good'' or ``bad''.}

% ---------------------------------------------------------------------
\subsection{Failure Modes}
\label{sec:results:failures}
% ---------------------------------------------------------------------
% Answers: RQ3. Operational definitions must be fixed BEFORE counting
% (see ANALYSIS_PLAN.md). Repository vocabulary
% (EXPERIMENT_STATUS_CANONICAL.md §0): single-class collapse,
% Reliable-heavy collapse (R >= ~91%), abstention-heavy collapse,
% technical failure. Describe behavior, not causes.

% TODO: AUTO-GENERATE FROM RESULTS
% TODO: AUTO-GENERATE FROM RESULTS
% Fix the operational thresholds in ANALYSIS_PLAN.md BEFORE filling this in.
% Repository vocabulary: docs/EXPERIMENT_STATUS_CANONICAL.md §0.
% Include sub-5,747 models (LLaVA-OV, Gemma 3 at A1@1000; InternVL3 at
% balanced-100) in a separate block, clearly labeled with their N.
\begin{table}[t]
  \centering
  \small
  \caption{Observed failure modes by model and condition, using pre-registered operational definitions. Behavioral descriptors only; no causal claims. \AUTOGEN{caption details}}
  \label{tab:failure_modes}
  \begin{tabular}{@{}l l l r l@{}}
    \toprule
    \textbf{Model} & \textbf{Cond.} & \textbf{Mode} & \textbf{N} & \textbf{Defining statistic} \\
    \midrule
    \AUTOGENROW{5}{Reliable-class collapse / abstention-heavy / parsing failure rows} \\
    \bottomrule
  \end{tabular}
\end{table}


\subsubsection{Reliable-Class Collapse}
\label{sec:results:failures:collapse}
% TODO: AUTO-GENERATE FROM RESULTS
\TODO{Models/conditions with $\geq$ threshold \R{} share; their alignment accuracy relative to the always-\R{} prior (0.889 full cohort; 0.928 on the A1@1{,}000 slice, which must not be compared directly with full-cohort numbers).}

\subsubsection{Abstention-Heavy Behavior}
\label{sec:results:failures:abstention}
% TODO: AUTO-GENERATE FROM RESULTS
\TODO{Models/conditions where \U{} dominates; report coverage and how few decided samples their alignment metrics rest on.}

\subsubsection{Parsing and Instruction-Following Failures}
\label{sec:results:failures:parsing}
% TODO: AUTO-GENERATE FROM RESULTS
\TODO{Parse-error and inference-error counts per run (all completed full-scale runs are recorded with zero; qualification-stage failures: InternVL3-8B unparsable outputs on balanced-100; Llama-3.2-Vision free-text answers; Ministral~3 control characters before the normalizer). Source: \texttt{outputs/diagnostics/model\_qualification/}, \texttt{docs/INTERNVL3\_QUALIFICATION.md}, \texttt{docs/MODEL\_EXPANSION\_CAMPAIGN.md}.}

\subsubsection{Qualitative Visual Failure Cases}
\label{sec:results:failures:visual}
% TODO: AUTO-GENERATE FROM RESULTS
\begin{figure}[t]
  \centering
  \FIGPLACEHOLDER{Grid of visual failure cases, e.g.\ detections whose verdict flips between A2 and A5 for several models; each tile shows the A2 overlay and A5 crop with the verdicts and the semantic label (if reviewed). Selection rule must be stated and must not be cherry-picked.}
  \caption{Representative visual failure cases. \AUTOGEN{selection rule and caption}}
  \label{fig:visual_failures}
\end{figure}

% ---------------------------------------------------------------------
\subsection{IoU-Based versus Semantic Evaluation}
\label{sec:results:iou_vs_semantic}
% ---------------------------------------------------------------------
% Answers: RQ4.
% TERMINOLOGY (mandatory):
%   "LabelMe-unmatched detection" = IoU-matching outcome only.
%   "semantic non-palm"           = human review label non_palm.
%   An unmatched detection is NOT a false positive unless the semantic
%   review labels it non_palm. Protocol-v2 "FP"/"TN"/"specificity" are
%   alignment quantities.

Two references are distinguished throughout this subsection. A
\emph{\unmatched{} detection} is one that IoU matching
(\cref{sec:method:iou}) did not pair with any LabelMe palm box; this is a
statement about annotation alignment only. A \emph{semantic non-palm} is a
detection that the human reviewer labeled \texttt{non\_palm}
(\cref{sec:method:semantic}). Alignment ``false positives'' (\R{} on a
\unmatched{} detection) are interpreted as detector false positives only to the
extent that the semantic review supports that interpretation.

% TODO: AUTO-GENERATE FROM RESULTS
% TODO: AUTO-GENERATE FROM RESULTS
% Sources (read-only; verify SHA-256 listed in docs/SEMANTIC_VALIDITY_AUDIT.md §G):
%   outputs/semantic_gt_review/human_review.csv            (638 unmatched)
%   outputs/semantic_gt_review/human_confidence_pilot.csv  (400 pilot)
% The pilot is outside the 5,747 cohort; report per-bin counts only
% (pooled shares are NOT prevalence because bins were sampled equally).
% Do NOT analyse the familiarization snapshot.
\begin{table}[t]
  \centering
  \small
  \caption{Human semantic review. Top: all \unmatched{} detections of the cohort. Bottom: lower-confidence pilot (outside the cohort; per-bin counts, not prevalence). \AUTOGEN{caption details}}
  \label{tab:semantic_review}
  \begin{tabular}{@{}l r r r r@{}}
    \toprule
    \textbf{Population} & \textbf{n} & \textbf{palm} & \textbf{ambiguous} & \textbf{non\_palm} \\
    \midrule
    \unmatched{} (conf.\ $\geq 0.5$) & \multicolumn{4}{c}{\AUTOGEN{counts}} \\
    \quad of which balanced-100 ``negatives'' & \multicolumn{4}{c}{\AUTOGEN{counts}} \\
    \quad of which A1@1{,}000 slice & \multicolumn{4}{c}{\AUTOGEN{counts}} \\
    \midrule
    Pilot $[0.10, 0.20)$ & \multicolumn{4}{c}{\AUTOGEN{counts}} \\
    Pilot $[0.20, 0.30)$ & \multicolumn{4}{c}{\AUTOGEN{counts}} \\
    Pilot $[0.30, 0.40)$ & \multicolumn{4}{c}{\AUTOGEN{counts}} \\
    Pilot $[0.40, 0.50)$ & \multicolumn{4}{c}{\AUTOGEN{counts}} \\
    \bottomrule
  \end{tabular}
\end{table}


% TODO: AUTO-GENERATE FROM RESULTS
% TODO: AUTO-GENERATE FROM RESULTS
% Join: human_review.csv (sample_id -> semantic label) x stored decisions
% (outputs/verification/<model>/<exp>/A*/sample_*.json or Protocol-v2
% evaluation CSV verification_label column). Only the 638 unmatched.
% A related, earlier file exists:
%   outputs/semantic_gt_evaluation/negative_pool_decisions_by_semantic.csv
% but that directory is SUPERSEDED — recompute from the official labels.
\begin{table}[t]
  \centering
  \small
  \caption{VLM verdicts on the 638 \unmatched{} detections, by human semantic label. An \Ur{} verdict on a semantic palm counts as an alignment ``TN'' under Protocol~v2. \AUTOGEN{caption details}}
  \label{tab:semantic_by_decision}
  \begin{tabular}{@{}l l r r r r r r@{}}
    \toprule
    & & \multicolumn{3}{c}{\textbf{semantic palm}} & \multicolumn{3}{c}{\textbf{ambiguous}} \\
    \cmidrule(lr){3-5}\cmidrule(lr){6-8}
    \textbf{Model} & \textbf{Cond.} & \R{} & \U{} & \Ur{} & \R{} & \U{} & \Ur{} \\
    \midrule
    \AUTOGENROW{8}{one row per model $\times$ condition; add a non\_palm block only if any exist} \\
    \bottomrule
  \end{tabular}
\end{table}


% TODO: AUTO-GENERATE FROM RESULTS
% TODO: AUTO-GENERATE FROM RESULTS
% Source: outputs/semantic_gt_evaluation/iou_sensitivity/iou_sensitivity_summary.csv
%   produced by scripts/analysis/iou_sensitivity_analysis.py
%   (Protocol-v2 matching, threshold varied; canonical threshold stays 0.5).
% VERIFY provenance before use: output lives under a SUPERSEDED directory;
% confirm the summary does not depend on superseded semantic labels, and
% recompute the semantic re-match column against the OFFICIAL human_review.csv.
\begin{table}[t]
  \centering
  \small
  \caption{Sensitivity of the \matched{}/\unmatched{} split to the IoU threshold (sensitivity analysis only; the canonical threshold is 0.5). \AUTOGEN{caption details}}
  \label{tab:iou_sensitivity}
  \begin{tabular}{@{}r r r r r@{}}
    \toprule
    \textbf{IoU thr.} & \textbf{Matched} & \textbf{Unmatched} & \textbf{Unmatched$\rightarrow$matched} & \textbf{Matched$\rightarrow$unmatched} \\
    \midrule
    \AUTOGENROW{5}{rows for thresholds 0.3--0.7; last two columns are changes relative to the canonical 0.5 split} \\
    \bottomrule
  \end{tabular}
\end{table}


% TODO: AUTO-GENERATE FROM RESULTS
\begin{figure}[t]
  \centering
  \FIGPLACEHOLDER{For each model $\times$ condition: VLM verdict distribution on the 638 \unmatched{} detections, split by semantic label (palm / ambiguous / non-palm).}
  \caption{VLM verdicts on \unmatched{} detections by human semantic label. \AUTOGEN{caption details}}
  \label{fig:semantic_comparison}
\end{figure}

% TODO: AUTO-GENERATE FROM RESULTS
% TODO: AUTO-GENERATE FROM RESULTS
% Two kinds of confusion matrix, kept visually separate:
%  (1) Alignment: verdict (R/U/Ur, 3 rows) x {LabelMe-matched, LabelMe-unmatched}
%      per model x condition. Source: Protocol-v2 A*_evaluation.csv.
%      Keep the Uncertain row (do not drop it) so coverage is visible.
%  (2) Semantic (unmatched subset only): verdict x {palm, ambiguous, non_palm}.
%      See tab_semantic_by_decision.tex.
% Where appropriate: only for models with substantive use of >= 2 verdicts;
% for collapsed models a confusion matrix is uninformative — cite
% tab_failure_modes instead. Consider moving the full set to the appendix.
\begin{table}[t]
  \centering
  \small
  \caption{Verdict $\times$ annotation-alignment label for selected model $\times$ condition cells. Columns are IoU-matching outcomes, not semantic classes. \AUTOGEN{caption details}}
  \label{tab:confusion}
  \begin{tabular}{@{}l l l r r@{}}
    \toprule
    \textbf{Model} & \textbf{Cond.} & \textbf{Verdict} & \textbf{\matched{}} & \textbf{\unmatched{}} \\
    \midrule
    \AUTOGENROW{5}{3 rows (R, U, Ur) per model $\times$ condition} \\
    \bottomrule
  \end{tabular}
\end{table}


\TODO{Report how many \unmatched{} detections are semantic palms, ambiguous, or non-palms; how the IoU threshold changes the matched/unmatched split; and what this implies for the meaning of alignment specificity. Do not compute semantic non-palm specificity if the cohort contains no confirmed non-palms.}

% ---------------------------------------------------------------------
\subsection{Qualitative Case Studies}
\label{sec:results:cases}
% ---------------------------------------------------------------------
% Selection must be rule-based and reproducible (record sample IDs and the
% rule in the appendix). Candidate rules:
%  (a) random sample (fixed seed) within strata of
%      {matched, unmatched-palm, unmatched-ambiguous} x {stable, variable};
%  (b) detections with maximal cross-model disagreement;
%  (c) detections whose verdict flips A2 -> A5 for most models.
% Show the A-condition images, every model's verdict, and (optionally) the
% model's visual_reasoning excerpt (quoted verbatim, not paraphrased).

% TODO: AUTO-GENERATE FROM RESULTS
\begin{figure}[t]
  \centering
  \FIGPLACEHOLDER{Case study 1: one detection shown under A1--A5 with all models' verdicts and verbatim reasoning excerpts.}
  \caption{Case study. \AUTOGEN{sample ID, selection rule, caption}}
  \label{fig:case1}
\end{figure}

% TODO: AUTO-GENERATE FROM RESULTS
\begin{figure}[t]
  \centering
  \FIGPLACEHOLDER{Case study 2: a LabelMe-unmatched detection reviewed as palm, with verdicts across models and conditions.}
  \caption{Case study. \AUTOGEN{sample ID, selection rule, caption}}
  \label{fig:case2}
\end{figure}

\TODO{Write 1--2 paragraphs per case. Quote model reasoning verbatim; do not infer model ``intent''.}
